\cotefile{136}
\scanpage{204}
Endomorphismes de l'anneau différentiel $\frac12$ simplicial $A_{\ast}$

(NB. on ne demande plus de respecter la structure de $S$-algèbre, $S=A_0^{\deleted{\unsure{*}}}$, i.e. on ne demande plus que $\varphi_0$ induise l'identité sur $A_0$).

Rappelons-nous qu'un complexe de De Rham peut se construire chaque fois qu'on a un anneau comm. $S$ avec idéal $S^+$ à puiss. divisées, et
$t\in S^+$, via une \inserted{$S$-algèbre} \deleted{algèbre} différentielle \inserted{$\wedge$-graduée à p.d.} $\frac12$ simpliciale $A_\ast(S,S^+,t)$, où
\[
A_n(S,S^+,t)=S\{X_0,\ldots,X_n\}[dX_0,\ldots,dX_n]\big/\bigl(\textstyle\sum X_i-t,\sum dX_i\bigr)_{\mathrm{pd}}
\]
Si on a un hom. d'anneaux à p.d., $S\xrightarrow{u} S'$ ($u(S^+)\subset S'{}^+\ldots$) et $\lambda'\in S'$ tel que
\deleted{$\unsure{\lambda t\in S^+|t\in S^+|\lambda\in S^+}$} $u(t)=\lambda t'$, on \deleted{en} déduit un hom. \inserted{de $S$-alg. $\wedge$-graduées différentielles à p.d.}
$A_\ast(u,\lambda): A_\ast(S,S^+,t)\longrightarrow \unsure{A'_\ast|A_\ast}(S',S'{}^+,t')$, qui en degré $n$ est l'hom.
\diagram{\begin{tikzcd}[ampersand replacement=\&,row sep=small]
S\{X_0,\ldots,X_n\}[dX_0,\ldots,dX_n]/\text{---}\arrow[d]\\
S'\{X_0,\ldots,X_n\}[dX_0,\ldots,dX_n]/\text{---}
\end{tikzcd}}
donné par \deleted{$\unsure{X_0\longleftrightarrow}$} $u:\deleted{\unsure{u}}S\longrightarrow S'$ sur les scalaires,
\[
\begin{aligned}
X_0&\longmapsto\lambda X_0\quad\bigl(\text{donc }\textstyle\sum X_0-t\longmapsto\sum\lambda X_0-u(t)=\sum\lambda X_{0}-\lambda t'\\[-2pt]
&\hspace{21em}=\lambda(\textstyle\sum X_0-t')\bigr)\\
dX_0&\longmapsto\lambda dX_0.
\end{aligned}
\]
Ainsi $A_\ast(S,S^+,t)$ devient un foncteur covariant \deleted{de} la catégorie des triples $(S,S^+,t)$, où les hom sont les couples $(u,\lambda)$ \inserted{comme ci-dessus} (composés par $(u',\lambda')(u,\lambda)=(u'u,u'(\lambda)\lambda')$).

\uline{NB} En plus de sa graduation, $A_n(S,S^+,t)$ admet une filtration \inserted{croissante} naturelle, où $\mathrm{Fil}^{\unsure{n|r}}A_n(S,S^+,t)$ est \unsure{le|les} sous
module des expressions $\sum F_{i_0\ldots i_p}\,dX_{i_0}\wedge\ldots\wedge dX_{i_p}$, avec $F_{i_0\ldots i_p}\in S\{X_0,\ldots,X_n\}$ \unsure{étant|elt|élément} de ``degré'' $\leq n-\unsure{p|r}$, ce sont des sous-$S$ modules (mais pas des idéaux), on
a $\mathrm{Fil}^n\cdot\mathrm{Fil}^m\subset\mathrm{Fil}^{n+m}$, $\mathrm{Fil}^0=S$ (scalaires), $\mathrm{Fil}^n=0$ si $n<0$. Cette filtration est respectée par les $A_\ast(u,\lambda)$.
% archive sheet number: 203
